\chapter{Dãy số, cấp số cộng và cấp số nhân}

\section{Dãy số}

\subsection{Dãy số hữu hạn}
\noindent\hspace*{1cm}- Mỗi hàm số $u:\{1;2;3;...;m\} \rightarrow \mathbb{R}(m \in \mathbb{N}^*)$ được gọi là một dãy số hữu hạn.\\[6pt]
\noindent\hspace*{1cm}- Dạng khai triển: $u_1, u_2, u_3,..., u_m$.\\[6pt]
\noindent\hspace*{1cm}- Số $u_1$ gọi là số hạng đầu và số $u_m$ gọi là số hạng cuối của dãy số.\\[20pt]
\noindent\textbf{Ví dụ:} Cho một dãy số:
\[
    u_1 = 1,
    u_2 = 4,
    u_3 = 9,
    u_4 = 12,
    u_5 = 17,
    u_6 = 20,
    u_7 = 23
\]
\noindent\hspace*{1.5cm}+ Đây là một dãy số hữu hạn có 7 phần tử\\[6pt]
\noindent\hspace*{1.5cm}+ Dạng khai triển của nó sẽ là:
\[
    1, 4, 9, 12, 17, 20, 23
\]
\noindent\hspace*{1.5cm}+ $u_1 = 1$ là số hạng đầu và $u_7 = 23$  là số hạng cuối.

\subsection{Dãy số vô hạn}
\noindent\hspace*{1cm}- Mỗi hàm số $u: \mathbb{N}^* \rightarrow \mathbb{R}$ được gọi là một dãy số vô hạn (gọi tắt là dãy số).\\[6pt]
\noindent\hspace*{1cm}- Dạng khai triển: $u_1, u_2, u_3,..., u_n,...$.\\[6pt]
\noindent\hspace*{1cm}- Dãy số đó còn được viết tắt là $(u_n)$.\\[6pt]
\noindent\hspace*{1cm}- Số $u_1$ gọi là số hạng thứ nhất (hay số hạng đầu),\\[6pt]
\noindent\hspace*{1.15cm} Số $u_2$ gọi là số hạng thứ hai,\\[6pt]
\noindent\hspace*{1.15cm} Số $u_3$ gọi là số hạng thứ ba,\\[6pt]
\noindent\hspace*{1.15cm} ...........,\\[6pt]
\noindent\hspace*{1.15cm} Số $u_n$ gọi là số hạng thứ $n$ (và là số hạng tổng quát của dãy số).\\[6pt]







